\documentclass[12pt,a4paper]{article}

\usepackage[margin=1in]{geometry}
\usepackage[T1]{fontenc}
\usepackage[utf8]{inputenc}
\usepackage{lmodern}
\usepackage{hyperref}
\usepackage{booktabs}
\usepackage{array}
\usepackage{enumitem}
\usepackage{amsmath}
\usepackage{xcolor}
\usepackage{setspace}
\usepackage{verbatim}

\hypersetup{
  colorlinks=true,
  linkcolor=blue!50!black,
  urlcolor=blue!50!black,
  citecolor=blue!50!black
}

\setlength{\emergencystretch}{2em}
\setlength{\parskip}{6pt}
\setlength{\parindent}{0pt}
\sloppy

\title{QuickBites Assignment 4: Sharding Design and Evaluation Report}
\author{Team ScaleOps}
\date{17 April 2026}

\begin{document}

\begin{titlepage}
\centering
\vspace*{1.8cm}
{\fontsize{22}{26}\selectfont \textbf{QuickBites}}\\[0.4cm]
{\Large Assignment 4 Report}\\[0.2cm]
{\large Sharding Design and Evaluation}\\[1.6cm]

\begin{tabular}{ll}
Course: & CS 432 Databases \\
Team: & ScaleOps \\
Submission Date: & 18 April 2026 \\
\end{tabular}

\vspace{1.4cm}

\begin{minipage}{0.84\textwidth}
  \centering
  \begin{tabular}{@{} L{5.0cm} L{7.0cm} @{}}
  \toprule
  \textbf{Team Members} & \textbf{Role} \\
  \midrule
  Bhavya Parmar & Subtask 4\\
  Surriya Gokul &  Subtask 1\\
  Sonagiri Thoshith Kautilya &  Subtask 1 \& Report\\
  Tanishq Chaudhari & Subtask 2 \& 3\\
  Mohit & Subtask 2 \& 3\\
  \bottomrule
  \end{tabular}
  \end{minipage}

\vspace{1.4cm}

\begin{minipage}{0.9\textwidth}
\textbf{Repository:} \href{https://github.com/mohit-codes-21/QuickBites}{https://github.com/mohit-codes-21/QuickBites}\\
\textbf{Video:} \href{https://drive.google.com/file/d/1e2TOupRxlVEyKsvpMLmblW_UvzhAfaY6/view?usp=sharing}{Project Demonstration Link}
\end{minipage}

\vfill
\end{titlepage}

\tableofcontents
\newpage

\section{Executive Summary}

This report documents the sharding design implemented for QuickBites in Assignment 4.
The system uses hash-based sharding on \texttt{customerID} with three shard nodes.
Customer-scoped tables are distributed across shards, while non-customer reference and
operational tables remain in the local monolith database. Migration and verification were
automated with dedicated scripts, and query routing was integrated into the Flask backend.
Validation includes migration integrity checks and API-level integration tests for point
routing and cross-shard range queries.

\section{Shard Key Chosen and Justification}

\subsection{Selected Shard Key}

The selected shard key is \texttt{customerID}.

\subsection{Justification Against Selection Criteria}

\textbf{High cardinality:}
\texttt{customerID} is unique per customer, so distinct key count scales with user growth.
This supports balanced partitioning under hash routing.

\textbf{Query alignment:}
Customer-facing workflows are keyed by customer context (cart, addresses, payments,
customer orders). Routing by \texttt{customerID} keeps these operations single-shard in
the common case.

\textbf{Stability:}
\texttt{customerID} is immutable after account creation. Immutable keys avoid row movement
caused by key updates.

\subsection{Skew Expectation}

For the design-time seed analysis (from \texttt{sql/SQL\_Dump.sql}),
customer IDs 1..10 map as: shard 0 = 3 rows (3,6,9),
shard 1 = 4 rows (1,4,7,10), shard 2 = 3 rows (2,5,8).
This indicates low initial skew. Under larger datasets,
modulo-hash distribution is expected to stabilize further.

\section{Partitioning Strategy Used and Why}

\subsection{Strategy}

QuickBites uses hash-based partitioning with:
\[
\texttt{shard\_id} = \texttt{customerID} \bmod 3
\]

\subsection{Rationale}

\begin{itemize}[leftmargin=*]
  \item Deterministic routing with constant-time lookup.
  \item No directory service required.
  \item Better write spread than simple range partitioning under sequential IDs.
  \item Co-location of customer-scoped entities by key.
\end{itemize}

\subsection{Alternatives Considered}

\begin{itemize}[leftmargin=*]
  \item Range partitioning was not chosen due to potential growth hotspots.
  \item Directory-based partitioning was not chosen due to added metadata complexity
  and lookup overhead for this assignment scope.
\end{itemize}

\section{How Query Routing Is Implemented in the Application}

\subsection{Routing Layer}

Routing is implemented using:

\begin{itemize}[leftmargin=*]
  \item \texttt{shard\_router.py}
  \item \texttt{ShardRouter.shard\_for\_customer(customer\_id)}
  \item \texttt{ShardRouter.connect\_for\_customer(customer\_id)}
  \item \texttt{ShardRouter.table\_name(logical\_table, shard\_id)}
  \item \texttt{app/app.py}: \texttt{\_get\_customer\_shard\_context(customer\_id)}
\end{itemize}

\subsection{Table Routing Scope}

Customer-key routed tables:
\texttt{Customer}, \texttt{Address}, \texttt{CartItem}, \texttt{Payment}, \texttt{Orders}.

Local-only tables:
\texttt{Member}, \texttt{DeliveryPartner}, \texttt{Restaurant}, \texttt{MenuItem},
\texttt{OrderItem}, \texttt{Delivery\_Assignments}, \texttt{OrderRating}, \texttt{MenuItemRating}.

\subsection{Implemented Query Patterns}

\textbf{Point routing:}
\texttt{GET /api/sharded/route/customer/<customer\_id>} returns resolved shard and table.

\textbf{Single-shard customer operations:}
Examples include
\texttt{GET /api/sharded/customers/<customer\_id>},
\texttt{GET /api/customer/orders}, and
\texttt{GET /api/customer/payments/last}.

\textbf{Range fan-out query:}
\texttt{GET /api/sharded/customers/range} with params
\texttt{start}, \texttt{end}, \texttt{limit}.
The backend queries each shard, merges rows, applies global sort, and returns a bounded
result set.

\subsection{Routing Validation Tests}

Integration tests in \texttt{tests/check.py} validate:

\begin{itemize}[leftmargin=*]
  \item Correct shard chosen for point query routing.
  \item Physical row presence on expected shard only.
  \item Correct multi-shard range fan-out behavior.
  \item Result equivalence with source \texttt{Customer} table for same bounds.
\end{itemize}

Latest execution status (\texttt{tests/check.py}): \texttt{Ran 2 tests ... OK}.
An additional display-oriented runner \texttt{tests/check\_show.py} prints full
records for point and range scenarios for demonstration screenshots.

\section{Which SQL Shard Tables Were Created and How Data Was Migrated}

\subsection{Physical Shard Tables Created}

For each shard ID \{0,1,2\}, the following tables are created:

\begin{itemize}[leftmargin=*]
  \item \texttt{shard\_<id>\_customer}
  \item \texttt{shard\_<id>\_address}
  \item \texttt{shard\_<id>\_cartitem}
  \item \texttt{shard\_<id>\_payment}
  \item \texttt{shard\_<id>\_orders}
\end{itemize}

Total shard tables: 15.

\subsection{Migration Workflow}

Migration is implemented in \texttt{shard\_admin.py}:

\begin{enumerate}[leftmargin=*]
  \item \texttt{setup}: create shard tables on all nodes.
  \item \texttt{migrate}: read rows from source \texttt{QB}, compute shard by
  \texttt{customerID \% 3}, and insert into destination shard tables.
  \item \texttt{verify}: validate counts, duplicate keys, and wrong-shard placement.
\end{enumerate}

\subsection{Migration Results}

Source: \url{logs/subtask2_shard_report.json}

Partition-integrity source: \url{logs/subtask2_partition_verification.json}

Overall result: \texttt{ok = true}

\textbf{Important note on counts.}
The table below reflects verification after resetting source data from
\texttt{sql/SQL\_Dump.sql} and rerunning shard migration.
Observed counts match the seed baseline exactly:
\texttt{Customer=10}, \texttt{Address=20}, \texttt{CartItem=12},
\texttt{Payment=12}, \texttt{Orders=12}.

\begin{center}
\begin{tabular}{>{\raggedright\arraybackslash}p{2.1cm} >{\centering\arraybackslash}p{1.8cm} >{\centering\arraybackslash}p{3.0cm} >{\centering\arraybackslash}p{1.2cm} >{\centering\arraybackslash}p{1.8cm} >{\centering\arraybackslash}p{1.4cm}}
\toprule
Table & Source & Shard Counts (0/1/2) & Delta & Duplicate Keys & Bad Route \\
\midrule
customer & 10 & 3 / 4 / 3 & 0 & 0 & 0 \\
address  & 20 & 6 / 8 / 6 & 0 & 0 & 0 \\
cartitem & 12 & 4 / 5 / 3 & 0 & 0 & 0 \\
payment  & 12 & 4 / 5 / 3 & 0 & 0 & 0 \\
orders   & 12 & 4 / 5 / 3 & 0 & 0 & 0 \\
\bottomrule
\end{tabular}
\end{center}

Conclusion: no loss, no duplication, and no wrong-shard placement for migrated sharded tables.

\section{Sharding Approach Used and How Shard Isolation Was Achieved}

\subsection{Approach Used}

This implementation uses multiple MySQL instances/databases (not Dockerized shards).

\begin{itemize}[leftmargin=*]
  \item Source monolith: \texttt{127.0.0.1:3306} (DB: \texttt{QB})
  \item Shard 0: \texttt{10.0.116.184:3307} (DB: \texttt{ScaleOps})
  \item Shard 1: \texttt{10.0.116.184:3308} (DB: \texttt{ScaleOps})
  \item Shard 2: \texttt{10.0.116.184:3309} (DB: \texttt{ScaleOps})
\end{itemize}

\subsection{Shard Isolation Mechanisms}

\begin{enumerate}[leftmargin=*]
  \item Connection isolation: each request is routed to an explicit shard connection.
  \item Physical table isolation: shard-local table namespaces prevent overlap.
  \item Deterministic ownership: each row maps to one shard using hash modulo.
  \item Verification isolation checks: migration verification detects duplicates and
  wrong placement across shards.
\end{enumerate}

\section{Results of the Scalability and Trade-offs Analysis}

\subsection{Distribution and Skew}

Using migrated counts:

\textbf{Customer skew coefficient}
\[
\text{skew}_{customer} = \frac{\max(3,4,3)}{(10/3)} = \frac{4}{3.33} \approx 1.20
\]

\textbf{Orders skew coefficient}
\[
\text{skew}_{orders} = \frac{\max(4,5,3)}{(12/3)} = \frac{5}{4} = 1.25
\]

Interpretation:
customer and order distributions are reasonably balanced under hash routing,
with mild skew expected from small sample sizes.

\subsection{Query Behavior Under Scale}

\begin{itemize}[leftmargin=*]
  \item Point customer queries scale efficiently since they hit one shard only.
  \item Range queries scale as fan-out operations and require merge-sort in app layer.
  \item A dedicated benchmark script (\texttt{app/benchmark\_indexing.py}) is available
  for index-impact analysis on customer/order/payment access patterns.
\end{itemize}

\subsection{Trade-offs}

\textbf{Advantages}
\begin{itemize}[leftmargin=*]
  \item Fast deterministic routing.
  \item Reduced contention for customer-scoped traffic.
  \item Clear operational mapping from key to shard.
\end{itemize}

\textbf{Costs}
\begin{itemize}[leftmargin=*]
  \item Fan-out overhead for range and global analytics queries.
  \item Hotspot risk from uneven customer activity.
  \item Additional operational overhead versus single-database deployment.
\end{itemize}

\section{SubTask 4: Horizontal Scaling and Trade-off Reflection}

\subsection{Horizontal vs Vertical Scaling}

In vertical scaling, a single database instance is upgraded with more CPU, memory,
or storage, but all traffic still converges to one node. In the implemented
horizontal approach, customer-scoped rows are distributed across multiple shard
nodes using deterministic hash routing on \texttt{customerID}. This allows growth
by adding shard capacity rather than only upgrading one server.

For QuickBites workloads, point queries (cart, addresses, payment history, customer
order lookups) typically touch one shard and therefore benefit from distributed
load. In contrast, range queries require scatter-gather across shards and add
application-side merge overhead.

In practice, this creates two clear query classes in our system:
single-shard customer queries with low coordination cost, and multi-shard
administrative/analytic queries with higher merge and sorting overhead.
As shard count increases, the first class scales well, while the second class
needs careful limit controls and query planning.

\subsection{Consistency}

Within a single shard, transaction behavior follows MySQL ACID semantics. However,
global consistency across all shards is weaker for multi-step workflows that involve
both shard-local and local-monolith tables.

Consistency can be stressed in cases such as:

\begin{itemize}[leftmargin=*]
  \item Partial commit paths across two independent connections.
  \item Temporary stale reads when one replica or path lags another source of truth.
  \item Cross-shard reporting windows where one shard is updated slightly earlier than others.
\end{itemize}

For assignment scope, this is acceptable because dominant customer operations are
single-shard, but it highlights that strict global serializability is not guaranteed
without additional coordination.

For this reason, our implementation follows an application-level consistency model:
customer-facing paths are optimized for shard-local correctness, while cross-shard
views are treated as eventually aligned snapshots over the queried shards.

\subsection{Availability}

Sharding improves overall service availability under partial load because unaffected
shards can continue serving their own key ranges. If one shard goes down, only
customers mapped to that shard lose access to shard-backed operations, while other
shards and local-only tables remain reachable.

Operationally, this trades complete outage risk for partial outage risk:

\begin{itemize}[leftmargin=*]
  \item Better aggregate uptime for non-impacted shards.
  \item Reduced blast radius compared with single-node failure.
  \item Need for shard-aware error handling and user-facing degraded-mode responses.
\end{itemize}

Therefore, availability is improved at system level but becomes key-dependent:
service quality for a request depends on whether the request's shard set is healthy.
This is an expected operational characteristic of shared-nothing sharded systems.

\subsection{Partition Tolerance (Simulated Shard Failure)}

Under a simulated network partition or shard-node outage, the router may fail to
connect to one or more shard ports. In the current implementation, affected requests
generally fail fast for the impacted shard path, while healthy shards continue to
operate.

For fan-out range queries, partition tolerance currently means either:

\begin{itemize}[leftmargin=*]
  \item returning an error when any required shard is unreachable, or
  \item (future improvement) returning partial results with explicit shard-miss metadata.
\end{itemize}

This behavior reflects a practical CAP trade-off: during partition events, the
system prioritizes correctness of returned data for reachable shards over pretending
full-cluster completeness.

For future hardening, this can be extended with explicit partial-result contracts,
retry/backoff for transient shard failures, and health-aware routing decisions
that separate temporary network faults from persistent shard outages.

\subsection{SubTask 4 Summary}

The implemented sharding design demonstrates clear horizontal scaling benefits for
customer-keyed traffic, with expected trade-offs in global consistency guarantees,
failure handling complexity, and fan-out query cost. These outcomes are consistent
with standard distributed-database behavior and align with the assignment's required
analysis dimensions.

Overall, SubTask 4 confirms that our architecture is effective for workload
locality around \texttt{customerID}, while also making the operational trade-offs
explicit for scale growth, fault handling, and cross-shard query patterns.

\section{Observations and Limitations}

\subsection{Observations}

\begin{itemize}[leftmargin=*]
  \item Query routing, migration, and verification are integrated end-to-end.
  \item Migration checks currently report full correctness for sharded table scope.
  \item Routing test cases demonstrate both required behaviors: correct point routing and
  correct multi-shard range fan-out.
\end{itemize}

\subsection{Limitations}

\begin{itemize}[leftmargin=*]
  \item Only direct \texttt{customerID}-routed tables are physically sharded.
  \item Local tables are still required for non-customer-centric workflows.
  \item No automated rebalancing/resharding policy is implemented.
  \item Fan-out query cost grows with shard count and result-set width.
\end{itemize}

\section{Conclusion}

QuickBites now has a working sharding implementation based on
\texttt{customerID} hash partitioning across three shard nodes.
The design includes deterministic routing, isolated shard tables,
automated migration, and correctness verification.
Measured results show balanced customer distribution with known
mild skew in order volume, and expected trade-offs between
single-shard efficiency and multi-shard fan-out costs.

\appendix
\section{Reproducibility Commands}

\begin{verbatim}
./.venv/bin/python shard_admin.py setup --drop-existing
./.venv/bin/python shard_admin.py migrate
./.venv/bin/python shard_admin.py verify
./.venv/bin/python verify_partition_integrity.py

./.venv/bin/python -m unittest tests/check.py -v
./.venv/bin/python tests/check_show.py
\end{verbatim}

\end{document}